\documentclass[../main.tex]{subfiles}
\usepackage{amssymb}
\graphicspath{{\subfix{../pic/}}}
\begin{document}

\section{Zadání}

Naprogramujte v ANSI C přenositelnou konzolovou aplikaci, která jako\newline
\hfill vstup načte z parametru na příkazové řádce matematickou funkci ve tvaru\newline
$y = f(x); \, x, y \in \mathbb{R}$, provede její analýzu a vytvoří soubor ve formátu PostScript  s grafem této funkce na zvoleném definičním oboru.
\vspace{0.5em}

Program se bude spouštět příkazem \textbf{graph.exe ⟨func⟩ ⟨out-file⟩ [⟨limits⟩]}.
Symbol ⟨func⟩ zastupuje zápis matematické funkce jedné reálné nezá-vislé proměnné (funkce ve více dimenzích program řešit nebude, nalezne-li během analýzy zápisu funkce více nezávisle proměnných než jednu, vypíše srozumitelné chybové hlášení a skončí). Závisle proměnná (zde y) je implicitní, a proto se levá strana rovnosti v zápisu nebude uvádět. Symbol ⟨out-file⟩ zastupuje jméno výstupního postscriptového souboru.

\vspace{0.5em}
Výsledkem práce programu bude soubor ve formátu PostScript, který bude zobrazovat graf zadané matematické funkce s vyznačenými souřadnými osami a (aspoň) význačnými hodnotami definičního oboru a oboru hodnot funkce.

\vspace{0.5em}
Pokud nebudou na příkazové řádce uvedeny alespoň dva argumenty, vypiš-te chybové hlášení a stručný návod k použití programu v angličtině podle běžných zvyklostí.

\subsection{Specifikace vstupu programu}

Vstupem programu jsou pouze parametry na příkazové řádce: Prvním (a nejdůležitějším) parametrem je matematická funkce, kterou má program zpracovat. S ohledem na to, že může její zápis
obsahovat mezery, napište program tak, aby akceptoval jak zápis funkce obklopený uvozovkami,
tak bez nich, tj. aby bylo možné program spouštět jak zadáním příkazu \textbf{graph.exe x+xˆ2}, tak
příkazem \textbf{graph.exe "x + xˆ2"}. Také použití mezer v zápisu \\
\makebox[\linewidth][s]{funkce musí být libovolné.  Je také možné, že uživatel vašeho programu} \\ z nepořádnosti někde v zápisu mezery použije a někde zase ne, třeba takto:
\textbf{"sin (x ˆ2)*cos(x / 2.0)"} – tato varianta zápisu je také akceptovatelná.

\newpage

\subsection{Způsob zápisu funkcí}
V zápisu matematické funkce, jejíž graf bude program generovat, se mohou vyskytnout tyto symboly: konstanty, proměnná, aritmetické operátory, funkce, závorky. Jiné symboly zápis obsahovat nesmí – pokud se v zápisu objeví, měl by program reagovat srozumitelným chybovým hlášením.
\vspace{0.5em}

Konstanty ( mohou být ve všech akceptovatelných formátech zápisu celého či reálného čísla
v jazyce C, tj. např. i \textbf{.5E-02} či \textbf{1E0}. S ohledem na to, že má program řešit jen 2D grafy, jediná
možná nezávislá proměnná, která se může v zápisu funkce vyskytovat je x.
\vspace{0.5em}

následující tabulka ukazuje veškeré povolené matematické operace:
\begin{table}[h!]
    \centering
    \begin{tabular}{|c|c|c|}
        \hline
        \textbf{Operace} & \textbf{Zápis operátoru} & \textbf{Příklad} \\ \hline
        unární mínus & $-$ & $-x$ \\ \hline
        sčítání & $+$ & $x + 2$ \\ \hline
        odčítání & $-$ & $x - 2$ \\ \hline
        násobení & $*$ & $2 * x$ \\ \hline
        dělení & $/$ & $x / 2$ \\ \hline
        umocnění & ˆ & xˆ2 \\ \hline
    \end{tabular}
    \caption{Povolené matematické operace}
\end{table}

Jiné operátory nejsou povolené. Priorita operátorů je dána matematickými pravidly, případně
upravena závorkami ( a ) – jiné druhy závorek se v zápisu funkce nesmí vyskytovat.

\subsection{Specifikace rozsahu zobrazení}
Třetí – \textbf{nepovinný} – parametr [⟨limits⟩] na příkazové řádce slouží k~předá-ní informací o rozsazích zobrazení grafu analyzované funkce, tedy o horní a~dolní mezi definičního oboru a oboru
hodnot funkce. Má tvar:
\vspace{0.5em}

{\centering ⟨xdol⟩:⟨xhor⟩:⟨ydol⟩:⟨yhor⟩ \par}

\vspace{0.5em}

Například příklad pro spuštění se zadanými rozsahy může vypadat takto: \textbf{graph.exe "sin(xˆ2)*cos(x)"\:sin2cos.ps 6.28:12.56:-1:1}. Není-li rozsah uveden (protože se jedná o nepovinný parametr), použijte implicitní nastavení,
které nechť je pro definiční obor i obor hodnot ⟨$-10$; $10$⟩.

\newpage
\subsection{Matematické funkce v zápisu}
Váš program by měl akceptovat v zápisu zobrazované funkce některé běžně používané matematické funkce. Funkce uvedené v následujícím výčtu program musí akceptovat, jinak nebude uznán
za řádně dokončený: absolutní hodnota \textbf{abs}; funkce eˆx \textbf{exp}, přirozený logaritmus \textbf{ln}, dekadický
logaritmus \textbf{log}; goniometrické funkce \textbf{sin}, \textbf{cos}, \textbf{tan}; cyklometrické funkce \textbf{asin}, \textbf{acos}, \textbf{atan}; hyperbolometrické funkce \textbf{sinh}, \textbf{cosh}, \textbf{tanh}.
\subsection{Specifikace výstupu programu}
Výstup programu bude směrován na konzoli pouze v případě chyby v~zadá-ní parametrů. Pokud
budou všechny parametry akceptovatelné a funkce k~zobrazení syntakticky správně zapsaná, nemusí program vypisovat na konzoli nic. Pokud program neobdrží požadovaný počet parametrů na
příkazovém řádku či pokud je funkce špatně zapsána, měl by vypsat srozumitelné chybové hlášení
v angličtině a skončit nenulovým návratovým kódem.
\begin{table}[h!]
    \centering
    \renewcommand{\arraystretch}{1.2} % Nastaví hustotu řádků
    \setlength{\tabcolsep}{4pt} % Zmenší horizontální mezery mezi sloupci
    \resizebox{1\textwidth}{!}{ % Zmenší celou tabulku na 90 % šířky stránky
    \begin{tabular}{|p{8cm}|c|}
        \hline
        \textbf{Popis chybového stavu} & \textbf{Návratová hodnota funkce int main()} \\ \hline
        Bez chyby/korektní ukončení činnosti programu. & 0 \\ \hline
        Nebyly předány správné parametry na příkazové řádce při spuštění programu. & 1 \\ \hline
        Řetězec předaný programu jako první parametr není akceptovatelným zápisem matematické funkce. & 2 \\ \hline
        Řetězec předaný programu jako druhý parametr nemůže být v hostitelském operačním systému názvem souboru nebo takový soubor nelze vytvořit či do něj ukládat informace. & 3 \\ \hline
        Uvedený nepovinný třetí parametr nelze dekódovat jako rozsahy zobrazení. & 4 \\ \hline
    \end{tabular}
    }
    \caption{Chybové stavy a návratové hodnoty funkce}
\end{table}
\vspace{0.5em}

Výsledkem činnosti programu by měl být soubor ve formátu PostScript, který bude obsahovat
vykreslený graf zadané funkce.
\end{document}